% ============================================================================
% ESQUEMA (banca): o ciclo de vida do classificador — laço ativo -> parada
% (3 gatilhos) -> liberação (3 condições + medida única no teste) ->
% produção (monitoramento de deriva em 3 camadas) -> retreino do zero ->
% volta. Inserido no APÊNDICE G (a7-parada-drift/texto.tex) por ordem do
% autor (2026-08-23), dentro de figure:
%   % ============================================================================
% ESQUEMA (banca): o ciclo de vida do classificador — laço ativo -> parada
% (3 gatilhos) -> liberação (3 condições + medida única no teste) ->
% produção (monitoramento de deriva em 3 camadas) -> retreino do zero ->
% volta. Inserido no APÊNDICE G (a7-parada-drift/texto.tex) por ordem do
% autor (2026-08-23), dentro de figure:
%   % ============================================================================
% ESQUEMA (banca): o ciclo de vida do classificador — laço ativo -> parada
% (3 gatilhos) -> liberação (3 condições + medida única no teste) ->
% produção (monitoramento de deriva em 3 camadas) -> retreino do zero ->
% volta. Inserido no APÊNDICE G (a7-parada-drift/texto.tex) por ordem do
% autor (2026-08-23), dentro de figure:
%   % ============================================================================
% ESQUEMA (banca): o ciclo de vida do classificador — laço ativo -> parada
% (3 gatilhos) -> liberação (3 condições + medida única no teste) ->
% produção (monitoramento de deriva em 3 camadas) -> retreino do zero ->
% volta. Inserido no APÊNDICE G (a7-parada-drift/texto.tex) por ordem do
% autor (2026-08-23), dentro de figure:
%   \input{3-metodo/esquemas-propostos/esq-ciclo-vida-modelo}
% AUTOCONTIDO: sem \ref, sem códigos de experimento, sem números além dos
% símbolos do próprio apêndice (L_antigo, L_novo). Requer amsmath (\text),
% que a tese já carrega.
% Os vãos entre colunas (1,2-1,4cm) NÃO comportam rótulos de seta em corpo
% 12 — a semântica vai nas caixas e na nota sob o ciclo (achado visual).
% Uso standalone (pré-visualização):
%   \documentclass[tikz,border=6pt,12pt]{standalone}
%   \usetikzlibrary{arrows.meta, positioning}
%   \begin{document}\input{esq-ciclo-vida-modelo.tex}\end{document}
% GEOMETRIA CALIBRADA PARA CORPO 12 (ppginf/book 12pt, textwidth 16cm).
% Requer apenas arrows.meta e positioning (já em packages.tex). Legível em
% P&B: produção = cinza claro; retreino = cinza médio; operação = branco.
% ============================================================================
\begin{tikzpicture}[
    x=1cm, y=1cm,
    caixa/.style={draw, rounded corners=2pt, align=center, text width=40mm,
                  minimum height=12mm, inner sep=3pt, font=\footnotesize},
    prod/.style={caixa, fill=gray!8},
    retre/.style={caixa, fill=gray!14},
    seta/.style={-{Stealth[length=2.6mm]}, thick},
    rot/.style={font=\scriptsize, align=center, inner sep=2pt}
]
    % ------------------ linha de cima: do laço à liberação -----------------
    \node[caixa] (laco) at (0,0)
        {\textbf{laço de aprendizado}\\\textbf{ativo}: seleção $\to$\\
         oráculo $\to$ retreino};
    \node[caixa] (parada) at (5.4,0)
        {\textbf{parada} pelo primeiro\\de três gatilhos:\\
         estagnação, orçamento\\ou exaustão do \textit{pool}};
    \node[caixa] (lib) at (10.8,0)
        {\textbf{liberação}: três condições;\\depois, uma única medida\\
         no teste intocado};

    % ------------------ linha de baixo: produção e retreino ----------------
    \node[prod] (producao) at (10.8,-3.6)
        {\textbf{produção}: monitora a\\deriva em três camadas\\
         (entrada; confiança;\\amostra com oráculo)};
    \node[retre, text width=44mm] (retreino) at (5.4,-3.6)
        {\textbf{retreino do zero} com\\
         \mbox{$L_{\text{antigo}} \cup L_{\text{novo}}$}, ciclo\\
         incremental no fluxo recente};

    % ------------------ setas do ciclo -------------------------------------
    \draw[seta] (laco) -- (parada);
    \draw[seta] (parada) -- (lib);
    \draw[seta] (lib) -- (producao);
    \draw[seta] (producao) -- (retreino);
    % o retorno fecha o ciclo pela esquerda, entrando por baixo do laço
    \draw[seta] (retreino.west) -| (laco.south);

    % ------------------ nota do ciclo de deriva ----------------------------
    \node[rot, text width=118mm] at (5.4,-5.6)
        {no monitoramento, as camadas 1 e 2 (custo zero) antecipam; a
         camada 3 (amostra com oráculo) confirma e dispara o retreino, que
         preserva todos os rótulos históricos e devolve o modelo à produção.};
\end{tikzpicture}

% AUTOCONTIDO: sem \ref, sem códigos de experimento, sem números além dos
% símbolos do próprio apêndice (L_antigo, L_novo). Requer amsmath (\text),
% que a tese já carrega.
% Os vãos entre colunas (1,2-1,4cm) NÃO comportam rótulos de seta em corpo
% 12 — a semântica vai nas caixas e na nota sob o ciclo (achado visual).
% Uso standalone (pré-visualização):
%   \documentclass[tikz,border=6pt,12pt]{standalone}
%   \usetikzlibrary{arrows.meta, positioning}
%   \begin{document}% ============================================================================
% ESQUEMA (banca): o ciclo de vida do classificador — laço ativo -> parada
% (3 gatilhos) -> liberação (3 condições + medida única no teste) ->
% produção (monitoramento de deriva em 3 camadas) -> retreino do zero ->
% volta. Inserido no APÊNDICE G (a7-parada-drift/texto.tex) por ordem do
% autor (2026-08-23), dentro de figure:
%   \input{3-metodo/esquemas-propostos/esq-ciclo-vida-modelo}
% AUTOCONTIDO: sem \ref, sem códigos de experimento, sem números além dos
% símbolos do próprio apêndice (L_antigo, L_novo). Requer amsmath (\text),
% que a tese já carrega.
% Os vãos entre colunas (1,2-1,4cm) NÃO comportam rótulos de seta em corpo
% 12 — a semântica vai nas caixas e na nota sob o ciclo (achado visual).
% Uso standalone (pré-visualização):
%   \documentclass[tikz,border=6pt,12pt]{standalone}
%   \usetikzlibrary{arrows.meta, positioning}
%   \begin{document}\input{esq-ciclo-vida-modelo.tex}\end{document}
% GEOMETRIA CALIBRADA PARA CORPO 12 (ppginf/book 12pt, textwidth 16cm).
% Requer apenas arrows.meta e positioning (já em packages.tex). Legível em
% P&B: produção = cinza claro; retreino = cinza médio; operação = branco.
% ============================================================================
\begin{tikzpicture}[
    x=1cm, y=1cm,
    caixa/.style={draw, rounded corners=2pt, align=center, text width=40mm,
                  minimum height=12mm, inner sep=3pt, font=\footnotesize},
    prod/.style={caixa, fill=gray!8},
    retre/.style={caixa, fill=gray!14},
    seta/.style={-{Stealth[length=2.6mm]}, thick},
    rot/.style={font=\scriptsize, align=center, inner sep=2pt}
]
    % ------------------ linha de cima: do laço à liberação -----------------
    \node[caixa] (laco) at (0,0)
        {\textbf{laço de aprendizado}\\\textbf{ativo}: seleção $\to$\\
         oráculo $\to$ retreino};
    \node[caixa] (parada) at (5.4,0)
        {\textbf{parada} pelo primeiro\\de três gatilhos:\\
         estagnação, orçamento\\ou exaustão do \textit{pool}};
    \node[caixa] (lib) at (10.8,0)
        {\textbf{liberação}: três condições;\\depois, uma única medida\\
         no teste intocado};

    % ------------------ linha de baixo: produção e retreino ----------------
    \node[prod] (producao) at (10.8,-3.6)
        {\textbf{produção}: monitora a\\deriva em três camadas\\
         (entrada; confiança;\\amostra com oráculo)};
    \node[retre, text width=44mm] (retreino) at (5.4,-3.6)
        {\textbf{retreino do zero} com\\
         \mbox{$L_{\text{antigo}} \cup L_{\text{novo}}$}, ciclo\\
         incremental no fluxo recente};

    % ------------------ setas do ciclo -------------------------------------
    \draw[seta] (laco) -- (parada);
    \draw[seta] (parada) -- (lib);
    \draw[seta] (lib) -- (producao);
    \draw[seta] (producao) -- (retreino);
    % o retorno fecha o ciclo pela esquerda, entrando por baixo do laço
    \draw[seta] (retreino.west) -| (laco.south);

    % ------------------ nota do ciclo de deriva ----------------------------
    \node[rot, text width=118mm] at (5.4,-5.6)
        {no monitoramento, as camadas 1 e 2 (custo zero) antecipam; a
         camada 3 (amostra com oráculo) confirma e dispara o retreino, que
         preserva todos os rótulos históricos e devolve o modelo à produção.};
\end{tikzpicture}
\end{document}
% GEOMETRIA CALIBRADA PARA CORPO 12 (ppginf/book 12pt, textwidth 16cm).
% Requer apenas arrows.meta e positioning (já em packages.tex). Legível em
% P&B: produção = cinza claro; retreino = cinza médio; operação = branco.
% ============================================================================
\begin{tikzpicture}[
    x=1cm, y=1cm,
    caixa/.style={draw, rounded corners=2pt, align=center, text width=40mm,
                  minimum height=12mm, inner sep=3pt, font=\footnotesize},
    prod/.style={caixa, fill=gray!8},
    retre/.style={caixa, fill=gray!14},
    seta/.style={-{Stealth[length=2.6mm]}, thick},
    rot/.style={font=\scriptsize, align=center, inner sep=2pt}
]
    % ------------------ linha de cima: do laço à liberação -----------------
    \node[caixa] (laco) at (0,0)
        {\textbf{laço de aprendizado}\\\textbf{ativo}: seleção $\to$\\
         oráculo $\to$ retreino};
    \node[caixa] (parada) at (5.4,0)
        {\textbf{parada} pelo primeiro\\de três gatilhos:\\
         estagnação, orçamento\\ou exaustão do \textit{pool}};
    \node[caixa] (lib) at (10.8,0)
        {\textbf{liberação}: três condições;\\depois, uma única medida\\
         no teste intocado};

    % ------------------ linha de baixo: produção e retreino ----------------
    \node[prod] (producao) at (10.8,-3.6)
        {\textbf{produção}: monitora a\\deriva em três camadas\\
         (entrada; confiança;\\amostra com oráculo)};
    \node[retre, text width=44mm] (retreino) at (5.4,-3.6)
        {\textbf{retreino do zero} com\\
         \mbox{$L_{\text{antigo}} \cup L_{\text{novo}}$}, ciclo\\
         incremental no fluxo recente};

    % ------------------ setas do ciclo -------------------------------------
    \draw[seta] (laco) -- (parada);
    \draw[seta] (parada) -- (lib);
    \draw[seta] (lib) -- (producao);
    \draw[seta] (producao) -- (retreino);
    % o retorno fecha o ciclo pela esquerda, entrando por baixo do laço
    \draw[seta] (retreino.west) -| (laco.south);

    % ------------------ nota do ciclo de deriva ----------------------------
    \node[rot, text width=118mm] at (5.4,-5.6)
        {no monitoramento, as camadas 1 e 2 (custo zero) antecipam; a
         camada 3 (amostra com oráculo) confirma e dispara o retreino, que
         preserva todos os rótulos históricos e devolve o modelo à produção.};
\end{tikzpicture}

% AUTOCONTIDO: sem \ref, sem códigos de experimento, sem números além dos
% símbolos do próprio apêndice (L_antigo, L_novo). Requer amsmath (\text),
% que a tese já carrega.
% Os vãos entre colunas (1,2-1,4cm) NÃO comportam rótulos de seta em corpo
% 12 — a semântica vai nas caixas e na nota sob o ciclo (achado visual).
% Uso standalone (pré-visualização):
%   \documentclass[tikz,border=6pt,12pt]{standalone}
%   \usetikzlibrary{arrows.meta, positioning}
%   \begin{document}% ============================================================================
% ESQUEMA (banca): o ciclo de vida do classificador — laço ativo -> parada
% (3 gatilhos) -> liberação (3 condições + medida única no teste) ->
% produção (monitoramento de deriva em 3 camadas) -> retreino do zero ->
% volta. Inserido no APÊNDICE G (a7-parada-drift/texto.tex) por ordem do
% autor (2026-08-23), dentro de figure:
%   % ============================================================================
% ESQUEMA (banca): o ciclo de vida do classificador — laço ativo -> parada
% (3 gatilhos) -> liberação (3 condições + medida única no teste) ->
% produção (monitoramento de deriva em 3 camadas) -> retreino do zero ->
% volta. Inserido no APÊNDICE G (a7-parada-drift/texto.tex) por ordem do
% autor (2026-08-23), dentro de figure:
%   \input{3-metodo/esquemas-propostos/esq-ciclo-vida-modelo}
% AUTOCONTIDO: sem \ref, sem códigos de experimento, sem números além dos
% símbolos do próprio apêndice (L_antigo, L_novo). Requer amsmath (\text),
% que a tese já carrega.
% Os vãos entre colunas (1,2-1,4cm) NÃO comportam rótulos de seta em corpo
% 12 — a semântica vai nas caixas e na nota sob o ciclo (achado visual).
% Uso standalone (pré-visualização):
%   \documentclass[tikz,border=6pt,12pt]{standalone}
%   \usetikzlibrary{arrows.meta, positioning}
%   \begin{document}\input{esq-ciclo-vida-modelo.tex}\end{document}
% GEOMETRIA CALIBRADA PARA CORPO 12 (ppginf/book 12pt, textwidth 16cm).
% Requer apenas arrows.meta e positioning (já em packages.tex). Legível em
% P&B: produção = cinza claro; retreino = cinza médio; operação = branco.
% ============================================================================
\begin{tikzpicture}[
    x=1cm, y=1cm,
    caixa/.style={draw, rounded corners=2pt, align=center, text width=40mm,
                  minimum height=12mm, inner sep=3pt, font=\footnotesize},
    prod/.style={caixa, fill=gray!8},
    retre/.style={caixa, fill=gray!14},
    seta/.style={-{Stealth[length=2.6mm]}, thick},
    rot/.style={font=\scriptsize, align=center, inner sep=2pt}
]
    % ------------------ linha de cima: do laço à liberação -----------------
    \node[caixa] (laco) at (0,0)
        {\textbf{laço de aprendizado}\\\textbf{ativo}: seleção $\to$\\
         oráculo $\to$ retreino};
    \node[caixa] (parada) at (5.4,0)
        {\textbf{parada} pelo primeiro\\de três gatilhos:\\
         estagnação, orçamento\\ou exaustão do \textit{pool}};
    \node[caixa] (lib) at (10.8,0)
        {\textbf{liberação}: três condições;\\depois, uma única medida\\
         no teste intocado};

    % ------------------ linha de baixo: produção e retreino ----------------
    \node[prod] (producao) at (10.8,-3.6)
        {\textbf{produção}: monitora a\\deriva em três camadas\\
         (entrada; confiança;\\amostra com oráculo)};
    \node[retre, text width=44mm] (retreino) at (5.4,-3.6)
        {\textbf{retreino do zero} com\\
         \mbox{$L_{\text{antigo}} \cup L_{\text{novo}}$}, ciclo\\
         incremental no fluxo recente};

    % ------------------ setas do ciclo -------------------------------------
    \draw[seta] (laco) -- (parada);
    \draw[seta] (parada) -- (lib);
    \draw[seta] (lib) -- (producao);
    \draw[seta] (producao) -- (retreino);
    % o retorno fecha o ciclo pela esquerda, entrando por baixo do laço
    \draw[seta] (retreino.west) -| (laco.south);

    % ------------------ nota do ciclo de deriva ----------------------------
    \node[rot, text width=118mm] at (5.4,-5.6)
        {no monitoramento, as camadas 1 e 2 (custo zero) antecipam; a
         camada 3 (amostra com oráculo) confirma e dispara o retreino, que
         preserva todos os rótulos históricos e devolve o modelo à produção.};
\end{tikzpicture}

% AUTOCONTIDO: sem \ref, sem códigos de experimento, sem números além dos
% símbolos do próprio apêndice (L_antigo, L_novo). Requer amsmath (\text),
% que a tese já carrega.
% Os vãos entre colunas (1,2-1,4cm) NÃO comportam rótulos de seta em corpo
% 12 — a semântica vai nas caixas e na nota sob o ciclo (achado visual).
% Uso standalone (pré-visualização):
%   \documentclass[tikz,border=6pt,12pt]{standalone}
%   \usetikzlibrary{arrows.meta, positioning}
%   \begin{document}% ============================================================================
% ESQUEMA (banca): o ciclo de vida do classificador — laço ativo -> parada
% (3 gatilhos) -> liberação (3 condições + medida única no teste) ->
% produção (monitoramento de deriva em 3 camadas) -> retreino do zero ->
% volta. Inserido no APÊNDICE G (a7-parada-drift/texto.tex) por ordem do
% autor (2026-08-23), dentro de figure:
%   \input{3-metodo/esquemas-propostos/esq-ciclo-vida-modelo}
% AUTOCONTIDO: sem \ref, sem códigos de experimento, sem números além dos
% símbolos do próprio apêndice (L_antigo, L_novo). Requer amsmath (\text),
% que a tese já carrega.
% Os vãos entre colunas (1,2-1,4cm) NÃO comportam rótulos de seta em corpo
% 12 — a semântica vai nas caixas e na nota sob o ciclo (achado visual).
% Uso standalone (pré-visualização):
%   \documentclass[tikz,border=6pt,12pt]{standalone}
%   \usetikzlibrary{arrows.meta, positioning}
%   \begin{document}\input{esq-ciclo-vida-modelo.tex}\end{document}
% GEOMETRIA CALIBRADA PARA CORPO 12 (ppginf/book 12pt, textwidth 16cm).
% Requer apenas arrows.meta e positioning (já em packages.tex). Legível em
% P&B: produção = cinza claro; retreino = cinza médio; operação = branco.
% ============================================================================
\begin{tikzpicture}[
    x=1cm, y=1cm,
    caixa/.style={draw, rounded corners=2pt, align=center, text width=40mm,
                  minimum height=12mm, inner sep=3pt, font=\footnotesize},
    prod/.style={caixa, fill=gray!8},
    retre/.style={caixa, fill=gray!14},
    seta/.style={-{Stealth[length=2.6mm]}, thick},
    rot/.style={font=\scriptsize, align=center, inner sep=2pt}
]
    % ------------------ linha de cima: do laço à liberação -----------------
    \node[caixa] (laco) at (0,0)
        {\textbf{laço de aprendizado}\\\textbf{ativo}: seleção $\to$\\
         oráculo $\to$ retreino};
    \node[caixa] (parada) at (5.4,0)
        {\textbf{parada} pelo primeiro\\de três gatilhos:\\
         estagnação, orçamento\\ou exaustão do \textit{pool}};
    \node[caixa] (lib) at (10.8,0)
        {\textbf{liberação}: três condições;\\depois, uma única medida\\
         no teste intocado};

    % ------------------ linha de baixo: produção e retreino ----------------
    \node[prod] (producao) at (10.8,-3.6)
        {\textbf{produção}: monitora a\\deriva em três camadas\\
         (entrada; confiança;\\amostra com oráculo)};
    \node[retre, text width=44mm] (retreino) at (5.4,-3.6)
        {\textbf{retreino do zero} com\\
         \mbox{$L_{\text{antigo}} \cup L_{\text{novo}}$}, ciclo\\
         incremental no fluxo recente};

    % ------------------ setas do ciclo -------------------------------------
    \draw[seta] (laco) -- (parada);
    \draw[seta] (parada) -- (lib);
    \draw[seta] (lib) -- (producao);
    \draw[seta] (producao) -- (retreino);
    % o retorno fecha o ciclo pela esquerda, entrando por baixo do laço
    \draw[seta] (retreino.west) -| (laco.south);

    % ------------------ nota do ciclo de deriva ----------------------------
    \node[rot, text width=118mm] at (5.4,-5.6)
        {no monitoramento, as camadas 1 e 2 (custo zero) antecipam; a
         camada 3 (amostra com oráculo) confirma e dispara o retreino, que
         preserva todos os rótulos históricos e devolve o modelo à produção.};
\end{tikzpicture}
\end{document}
% GEOMETRIA CALIBRADA PARA CORPO 12 (ppginf/book 12pt, textwidth 16cm).
% Requer apenas arrows.meta e positioning (já em packages.tex). Legível em
% P&B: produção = cinza claro; retreino = cinza médio; operação = branco.
% ============================================================================
\begin{tikzpicture}[
    x=1cm, y=1cm,
    caixa/.style={draw, rounded corners=2pt, align=center, text width=40mm,
                  minimum height=12mm, inner sep=3pt, font=\footnotesize},
    prod/.style={caixa, fill=gray!8},
    retre/.style={caixa, fill=gray!14},
    seta/.style={-{Stealth[length=2.6mm]}, thick},
    rot/.style={font=\scriptsize, align=center, inner sep=2pt}
]
    % ------------------ linha de cima: do laço à liberação -----------------
    \node[caixa] (laco) at (0,0)
        {\textbf{laço de aprendizado}\\\textbf{ativo}: seleção $\to$\\
         oráculo $\to$ retreino};
    \node[caixa] (parada) at (5.4,0)
        {\textbf{parada} pelo primeiro\\de três gatilhos:\\
         estagnação, orçamento\\ou exaustão do \textit{pool}};
    \node[caixa] (lib) at (10.8,0)
        {\textbf{liberação}: três condições;\\depois, uma única medida\\
         no teste intocado};

    % ------------------ linha de baixo: produção e retreino ----------------
    \node[prod] (producao) at (10.8,-3.6)
        {\textbf{produção}: monitora a\\deriva em três camadas\\
         (entrada; confiança;\\amostra com oráculo)};
    \node[retre, text width=44mm] (retreino) at (5.4,-3.6)
        {\textbf{retreino do zero} com\\
         \mbox{$L_{\text{antigo}} \cup L_{\text{novo}}$}, ciclo\\
         incremental no fluxo recente};

    % ------------------ setas do ciclo -------------------------------------
    \draw[seta] (laco) -- (parada);
    \draw[seta] (parada) -- (lib);
    \draw[seta] (lib) -- (producao);
    \draw[seta] (producao) -- (retreino);
    % o retorno fecha o ciclo pela esquerda, entrando por baixo do laço
    \draw[seta] (retreino.west) -| (laco.south);

    % ------------------ nota do ciclo de deriva ----------------------------
    \node[rot, text width=118mm] at (5.4,-5.6)
        {no monitoramento, as camadas 1 e 2 (custo zero) antecipam; a
         camada 3 (amostra com oráculo) confirma e dispara o retreino, que
         preserva todos os rótulos históricos e devolve o modelo à produção.};
\end{tikzpicture}
\end{document}
% GEOMETRIA CALIBRADA PARA CORPO 12 (ppginf/book 12pt, textwidth 16cm).
% Requer apenas arrows.meta e positioning (já em packages.tex). Legível em
% P&B: produção = cinza claro; retreino = cinza médio; operação = branco.
% ============================================================================
\begin{tikzpicture}[
    x=1cm, y=1cm,
    caixa/.style={draw, rounded corners=2pt, align=center, text width=40mm,
                  minimum height=12mm, inner sep=3pt, font=\footnotesize},
    prod/.style={caixa, fill=gray!8},
    retre/.style={caixa, fill=gray!14},
    seta/.style={-{Stealth[length=2.6mm]}, thick},
    rot/.style={font=\scriptsize, align=center, inner sep=2pt}
]
    % ------------------ linha de cima: do laço à liberação -----------------
    \node[caixa] (laco) at (0,0)
        {\textbf{laço de aprendizado}\\\textbf{ativo}: seleção $\to$\\
         oráculo $\to$ retreino};
    \node[caixa] (parada) at (5.4,0)
        {\textbf{parada} pelo primeiro\\de três gatilhos:\\
         estagnação, orçamento\\ou exaustão do \textit{pool}};
    \node[caixa] (lib) at (10.8,0)
        {\textbf{liberação}: três condições;\\depois, uma única medida\\
         no teste intocado};

    % ------------------ linha de baixo: produção e retreino ----------------
    \node[prod] (producao) at (10.8,-3.6)
        {\textbf{produção}: monitora a\\deriva em três camadas\\
         (entrada; confiança;\\amostra com oráculo)};
    \node[retre, text width=44mm] (retreino) at (5.4,-3.6)
        {\textbf{retreino do zero} com\\
         \mbox{$L_{\text{antigo}} \cup L_{\text{novo}}$}, ciclo\\
         incremental no fluxo recente};

    % ------------------ setas do ciclo -------------------------------------
    \draw[seta] (laco) -- (parada);
    \draw[seta] (parada) -- (lib);
    \draw[seta] (lib) -- (producao);
    \draw[seta] (producao) -- (retreino);
    % o retorno fecha o ciclo pela esquerda, entrando por baixo do laço
    \draw[seta] (retreino.west) -| (laco.south);

    % ------------------ nota do ciclo de deriva ----------------------------
    \node[rot, text width=118mm] at (5.4,-5.6)
        {no monitoramento, as camadas 1 e 2 (custo zero) antecipam; a
         camada 3 (amostra com oráculo) confirma e dispara o retreino, que
         preserva todos os rótulos históricos e devolve o modelo à produção.};
\end{tikzpicture}

% AUTOCONTIDO: sem \ref, sem códigos de experimento, sem números além dos
% símbolos do próprio apêndice (L_antigo, L_novo). Requer amsmath (\text),
% que a tese já carrega.
% Os vãos entre colunas (1,2-1,4cm) NÃO comportam rótulos de seta em corpo
% 12 — a semântica vai nas caixas e na nota sob o ciclo (achado visual).
% Uso standalone (pré-visualização):
%   \documentclass[tikz,border=6pt,12pt]{standalone}
%   \usetikzlibrary{arrows.meta, positioning}
%   \begin{document}% ============================================================================
% ESQUEMA (banca): o ciclo de vida do classificador — laço ativo -> parada
% (3 gatilhos) -> liberação (3 condições + medida única no teste) ->
% produção (monitoramento de deriva em 3 camadas) -> retreino do zero ->
% volta. Inserido no APÊNDICE G (a7-parada-drift/texto.tex) por ordem do
% autor (2026-08-23), dentro de figure:
%   % ============================================================================
% ESQUEMA (banca): o ciclo de vida do classificador — laço ativo -> parada
% (3 gatilhos) -> liberação (3 condições + medida única no teste) ->
% produção (monitoramento de deriva em 3 camadas) -> retreino do zero ->
% volta. Inserido no APÊNDICE G (a7-parada-drift/texto.tex) por ordem do
% autor (2026-08-23), dentro de figure:
%   % ============================================================================
% ESQUEMA (banca): o ciclo de vida do classificador — laço ativo -> parada
% (3 gatilhos) -> liberação (3 condições + medida única no teste) ->
% produção (monitoramento de deriva em 3 camadas) -> retreino do zero ->
% volta. Inserido no APÊNDICE G (a7-parada-drift/texto.tex) por ordem do
% autor (2026-08-23), dentro de figure:
%   \input{3-metodo/esquemas-propostos/esq-ciclo-vida-modelo}
% AUTOCONTIDO: sem \ref, sem códigos de experimento, sem números além dos
% símbolos do próprio apêndice (L_antigo, L_novo). Requer amsmath (\text),
% que a tese já carrega.
% Os vãos entre colunas (1,2-1,4cm) NÃO comportam rótulos de seta em corpo
% 12 — a semântica vai nas caixas e na nota sob o ciclo (achado visual).
% Uso standalone (pré-visualização):
%   \documentclass[tikz,border=6pt,12pt]{standalone}
%   \usetikzlibrary{arrows.meta, positioning}
%   \begin{document}\input{esq-ciclo-vida-modelo.tex}\end{document}
% GEOMETRIA CALIBRADA PARA CORPO 12 (ppginf/book 12pt, textwidth 16cm).
% Requer apenas arrows.meta e positioning (já em packages.tex). Legível em
% P&B: produção = cinza claro; retreino = cinza médio; operação = branco.
% ============================================================================
\begin{tikzpicture}[
    x=1cm, y=1cm,
    caixa/.style={draw, rounded corners=2pt, align=center, text width=40mm,
                  minimum height=12mm, inner sep=3pt, font=\footnotesize},
    prod/.style={caixa, fill=gray!8},
    retre/.style={caixa, fill=gray!14},
    seta/.style={-{Stealth[length=2.6mm]}, thick},
    rot/.style={font=\scriptsize, align=center, inner sep=2pt}
]
    % ------------------ linha de cima: do laço à liberação -----------------
    \node[caixa] (laco) at (0,0)
        {\textbf{laço de aprendizado}\\\textbf{ativo}: seleção $\to$\\
         oráculo $\to$ retreino};
    \node[caixa] (parada) at (5.4,0)
        {\textbf{parada} pelo primeiro\\de três gatilhos:\\
         estagnação, orçamento\\ou exaustão do \textit{pool}};
    \node[caixa] (lib) at (10.8,0)
        {\textbf{liberação}: três condições;\\depois, uma única medida\\
         no teste intocado};

    % ------------------ linha de baixo: produção e retreino ----------------
    \node[prod] (producao) at (10.8,-3.6)
        {\textbf{produção}: monitora a\\deriva em três camadas\\
         (entrada; confiança;\\amostra com oráculo)};
    \node[retre, text width=44mm] (retreino) at (5.4,-3.6)
        {\textbf{retreino do zero} com\\
         \mbox{$L_{\text{antigo}} \cup L_{\text{novo}}$}, ciclo\\
         incremental no fluxo recente};

    % ------------------ setas do ciclo -------------------------------------
    \draw[seta] (laco) -- (parada);
    \draw[seta] (parada) -- (lib);
    \draw[seta] (lib) -- (producao);
    \draw[seta] (producao) -- (retreino);
    % o retorno fecha o ciclo pela esquerda, entrando por baixo do laço
    \draw[seta] (retreino.west) -| (laco.south);

    % ------------------ nota do ciclo de deriva ----------------------------
    \node[rot, text width=118mm] at (5.4,-5.6)
        {no monitoramento, as camadas 1 e 2 (custo zero) antecipam; a
         camada 3 (amostra com oráculo) confirma e dispara o retreino, que
         preserva todos os rótulos históricos e devolve o modelo à produção.};
\end{tikzpicture}

% AUTOCONTIDO: sem \ref, sem códigos de experimento, sem números além dos
% símbolos do próprio apêndice (L_antigo, L_novo). Requer amsmath (\text),
% que a tese já carrega.
% Os vãos entre colunas (1,2-1,4cm) NÃO comportam rótulos de seta em corpo
% 12 — a semântica vai nas caixas e na nota sob o ciclo (achado visual).
% Uso standalone (pré-visualização):
%   \documentclass[tikz,border=6pt,12pt]{standalone}
%   \usetikzlibrary{arrows.meta, positioning}
%   \begin{document}% ============================================================================
% ESQUEMA (banca): o ciclo de vida do classificador — laço ativo -> parada
% (3 gatilhos) -> liberação (3 condições + medida única no teste) ->
% produção (monitoramento de deriva em 3 camadas) -> retreino do zero ->
% volta. Inserido no APÊNDICE G (a7-parada-drift/texto.tex) por ordem do
% autor (2026-08-23), dentro de figure:
%   \input{3-metodo/esquemas-propostos/esq-ciclo-vida-modelo}
% AUTOCONTIDO: sem \ref, sem códigos de experimento, sem números além dos
% símbolos do próprio apêndice (L_antigo, L_novo). Requer amsmath (\text),
% que a tese já carrega.
% Os vãos entre colunas (1,2-1,4cm) NÃO comportam rótulos de seta em corpo
% 12 — a semântica vai nas caixas e na nota sob o ciclo (achado visual).
% Uso standalone (pré-visualização):
%   \documentclass[tikz,border=6pt,12pt]{standalone}
%   \usetikzlibrary{arrows.meta, positioning}
%   \begin{document}\input{esq-ciclo-vida-modelo.tex}\end{document}
% GEOMETRIA CALIBRADA PARA CORPO 12 (ppginf/book 12pt, textwidth 16cm).
% Requer apenas arrows.meta e positioning (já em packages.tex). Legível em
% P&B: produção = cinza claro; retreino = cinza médio; operação = branco.
% ============================================================================
\begin{tikzpicture}[
    x=1cm, y=1cm,
    caixa/.style={draw, rounded corners=2pt, align=center, text width=40mm,
                  minimum height=12mm, inner sep=3pt, font=\footnotesize},
    prod/.style={caixa, fill=gray!8},
    retre/.style={caixa, fill=gray!14},
    seta/.style={-{Stealth[length=2.6mm]}, thick},
    rot/.style={font=\scriptsize, align=center, inner sep=2pt}
]
    % ------------------ linha de cima: do laço à liberação -----------------
    \node[caixa] (laco) at (0,0)
        {\textbf{laço de aprendizado}\\\textbf{ativo}: seleção $\to$\\
         oráculo $\to$ retreino};
    \node[caixa] (parada) at (5.4,0)
        {\textbf{parada} pelo primeiro\\de três gatilhos:\\
         estagnação, orçamento\\ou exaustão do \textit{pool}};
    \node[caixa] (lib) at (10.8,0)
        {\textbf{liberação}: três condições;\\depois, uma única medida\\
         no teste intocado};

    % ------------------ linha de baixo: produção e retreino ----------------
    \node[prod] (producao) at (10.8,-3.6)
        {\textbf{produção}: monitora a\\deriva em três camadas\\
         (entrada; confiança;\\amostra com oráculo)};
    \node[retre, text width=44mm] (retreino) at (5.4,-3.6)
        {\textbf{retreino do zero} com\\
         \mbox{$L_{\text{antigo}} \cup L_{\text{novo}}$}, ciclo\\
         incremental no fluxo recente};

    % ------------------ setas do ciclo -------------------------------------
    \draw[seta] (laco) -- (parada);
    \draw[seta] (parada) -- (lib);
    \draw[seta] (lib) -- (producao);
    \draw[seta] (producao) -- (retreino);
    % o retorno fecha o ciclo pela esquerda, entrando por baixo do laço
    \draw[seta] (retreino.west) -| (laco.south);

    % ------------------ nota do ciclo de deriva ----------------------------
    \node[rot, text width=118mm] at (5.4,-5.6)
        {no monitoramento, as camadas 1 e 2 (custo zero) antecipam; a
         camada 3 (amostra com oráculo) confirma e dispara o retreino, que
         preserva todos os rótulos históricos e devolve o modelo à produção.};
\end{tikzpicture}
\end{document}
% GEOMETRIA CALIBRADA PARA CORPO 12 (ppginf/book 12pt, textwidth 16cm).
% Requer apenas arrows.meta e positioning (já em packages.tex). Legível em
% P&B: produção = cinza claro; retreino = cinza médio; operação = branco.
% ============================================================================
\begin{tikzpicture}[
    x=1cm, y=1cm,
    caixa/.style={draw, rounded corners=2pt, align=center, text width=40mm,
                  minimum height=12mm, inner sep=3pt, font=\footnotesize},
    prod/.style={caixa, fill=gray!8},
    retre/.style={caixa, fill=gray!14},
    seta/.style={-{Stealth[length=2.6mm]}, thick},
    rot/.style={font=\scriptsize, align=center, inner sep=2pt}
]
    % ------------------ linha de cima: do laço à liberação -----------------
    \node[caixa] (laco) at (0,0)
        {\textbf{laço de aprendizado}\\\textbf{ativo}: seleção $\to$\\
         oráculo $\to$ retreino};
    \node[caixa] (parada) at (5.4,0)
        {\textbf{parada} pelo primeiro\\de três gatilhos:\\
         estagnação, orçamento\\ou exaustão do \textit{pool}};
    \node[caixa] (lib) at (10.8,0)
        {\textbf{liberação}: três condições;\\depois, uma única medida\\
         no teste intocado};

    % ------------------ linha de baixo: produção e retreino ----------------
    \node[prod] (producao) at (10.8,-3.6)
        {\textbf{produção}: monitora a\\deriva em três camadas\\
         (entrada; confiança;\\amostra com oráculo)};
    \node[retre, text width=44mm] (retreino) at (5.4,-3.6)
        {\textbf{retreino do zero} com\\
         \mbox{$L_{\text{antigo}} \cup L_{\text{novo}}$}, ciclo\\
         incremental no fluxo recente};

    % ------------------ setas do ciclo -------------------------------------
    \draw[seta] (laco) -- (parada);
    \draw[seta] (parada) -- (lib);
    \draw[seta] (lib) -- (producao);
    \draw[seta] (producao) -- (retreino);
    % o retorno fecha o ciclo pela esquerda, entrando por baixo do laço
    \draw[seta] (retreino.west) -| (laco.south);

    % ------------------ nota do ciclo de deriva ----------------------------
    \node[rot, text width=118mm] at (5.4,-5.6)
        {no monitoramento, as camadas 1 e 2 (custo zero) antecipam; a
         camada 3 (amostra com oráculo) confirma e dispara o retreino, que
         preserva todos os rótulos históricos e devolve o modelo à produção.};
\end{tikzpicture}

% AUTOCONTIDO: sem \ref, sem códigos de experimento, sem números além dos
% símbolos do próprio apêndice (L_antigo, L_novo). Requer amsmath (\text),
% que a tese já carrega.
% Os vãos entre colunas (1,2-1,4cm) NÃO comportam rótulos de seta em corpo
% 12 — a semântica vai nas caixas e na nota sob o ciclo (achado visual).
% Uso standalone (pré-visualização):
%   \documentclass[tikz,border=6pt,12pt]{standalone}
%   \usetikzlibrary{arrows.meta, positioning}
%   \begin{document}% ============================================================================
% ESQUEMA (banca): o ciclo de vida do classificador — laço ativo -> parada
% (3 gatilhos) -> liberação (3 condições + medida única no teste) ->
% produção (monitoramento de deriva em 3 camadas) -> retreino do zero ->
% volta. Inserido no APÊNDICE G (a7-parada-drift/texto.tex) por ordem do
% autor (2026-08-23), dentro de figure:
%   % ============================================================================
% ESQUEMA (banca): o ciclo de vida do classificador — laço ativo -> parada
% (3 gatilhos) -> liberação (3 condições + medida única no teste) ->
% produção (monitoramento de deriva em 3 camadas) -> retreino do zero ->
% volta. Inserido no APÊNDICE G (a7-parada-drift/texto.tex) por ordem do
% autor (2026-08-23), dentro de figure:
%   \input{3-metodo/esquemas-propostos/esq-ciclo-vida-modelo}
% AUTOCONTIDO: sem \ref, sem códigos de experimento, sem números além dos
% símbolos do próprio apêndice (L_antigo, L_novo). Requer amsmath (\text),
% que a tese já carrega.
% Os vãos entre colunas (1,2-1,4cm) NÃO comportam rótulos de seta em corpo
% 12 — a semântica vai nas caixas e na nota sob o ciclo (achado visual).
% Uso standalone (pré-visualização):
%   \documentclass[tikz,border=6pt,12pt]{standalone}
%   \usetikzlibrary{arrows.meta, positioning}
%   \begin{document}\input{esq-ciclo-vida-modelo.tex}\end{document}
% GEOMETRIA CALIBRADA PARA CORPO 12 (ppginf/book 12pt, textwidth 16cm).
% Requer apenas arrows.meta e positioning (já em packages.tex). Legível em
% P&B: produção = cinza claro; retreino = cinza médio; operação = branco.
% ============================================================================
\begin{tikzpicture}[
    x=1cm, y=1cm,
    caixa/.style={draw, rounded corners=2pt, align=center, text width=40mm,
                  minimum height=12mm, inner sep=3pt, font=\footnotesize},
    prod/.style={caixa, fill=gray!8},
    retre/.style={caixa, fill=gray!14},
    seta/.style={-{Stealth[length=2.6mm]}, thick},
    rot/.style={font=\scriptsize, align=center, inner sep=2pt}
]
    % ------------------ linha de cima: do laço à liberação -----------------
    \node[caixa] (laco) at (0,0)
        {\textbf{laço de aprendizado}\\\textbf{ativo}: seleção $\to$\\
         oráculo $\to$ retreino};
    \node[caixa] (parada) at (5.4,0)
        {\textbf{parada} pelo primeiro\\de três gatilhos:\\
         estagnação, orçamento\\ou exaustão do \textit{pool}};
    \node[caixa] (lib) at (10.8,0)
        {\textbf{liberação}: três condições;\\depois, uma única medida\\
         no teste intocado};

    % ------------------ linha de baixo: produção e retreino ----------------
    \node[prod] (producao) at (10.8,-3.6)
        {\textbf{produção}: monitora a\\deriva em três camadas\\
         (entrada; confiança;\\amostra com oráculo)};
    \node[retre, text width=44mm] (retreino) at (5.4,-3.6)
        {\textbf{retreino do zero} com\\
         \mbox{$L_{\text{antigo}} \cup L_{\text{novo}}$}, ciclo\\
         incremental no fluxo recente};

    % ------------------ setas do ciclo -------------------------------------
    \draw[seta] (laco) -- (parada);
    \draw[seta] (parada) -- (lib);
    \draw[seta] (lib) -- (producao);
    \draw[seta] (producao) -- (retreino);
    % o retorno fecha o ciclo pela esquerda, entrando por baixo do laço
    \draw[seta] (retreino.west) -| (laco.south);

    % ------------------ nota do ciclo de deriva ----------------------------
    \node[rot, text width=118mm] at (5.4,-5.6)
        {no monitoramento, as camadas 1 e 2 (custo zero) antecipam; a
         camada 3 (amostra com oráculo) confirma e dispara o retreino, que
         preserva todos os rótulos históricos e devolve o modelo à produção.};
\end{tikzpicture}

% AUTOCONTIDO: sem \ref, sem códigos de experimento, sem números além dos
% símbolos do próprio apêndice (L_antigo, L_novo). Requer amsmath (\text),
% que a tese já carrega.
% Os vãos entre colunas (1,2-1,4cm) NÃO comportam rótulos de seta em corpo
% 12 — a semântica vai nas caixas e na nota sob o ciclo (achado visual).
% Uso standalone (pré-visualização):
%   \documentclass[tikz,border=6pt,12pt]{standalone}
%   \usetikzlibrary{arrows.meta, positioning}
%   \begin{document}% ============================================================================
% ESQUEMA (banca): o ciclo de vida do classificador — laço ativo -> parada
% (3 gatilhos) -> liberação (3 condições + medida única no teste) ->
% produção (monitoramento de deriva em 3 camadas) -> retreino do zero ->
% volta. Inserido no APÊNDICE G (a7-parada-drift/texto.tex) por ordem do
% autor (2026-08-23), dentro de figure:
%   \input{3-metodo/esquemas-propostos/esq-ciclo-vida-modelo}
% AUTOCONTIDO: sem \ref, sem códigos de experimento, sem números além dos
% símbolos do próprio apêndice (L_antigo, L_novo). Requer amsmath (\text),
% que a tese já carrega.
% Os vãos entre colunas (1,2-1,4cm) NÃO comportam rótulos de seta em corpo
% 12 — a semântica vai nas caixas e na nota sob o ciclo (achado visual).
% Uso standalone (pré-visualização):
%   \documentclass[tikz,border=6pt,12pt]{standalone}
%   \usetikzlibrary{arrows.meta, positioning}
%   \begin{document}\input{esq-ciclo-vida-modelo.tex}\end{document}
% GEOMETRIA CALIBRADA PARA CORPO 12 (ppginf/book 12pt, textwidth 16cm).
% Requer apenas arrows.meta e positioning (já em packages.tex). Legível em
% P&B: produção = cinza claro; retreino = cinza médio; operação = branco.
% ============================================================================
\begin{tikzpicture}[
    x=1cm, y=1cm,
    caixa/.style={draw, rounded corners=2pt, align=center, text width=40mm,
                  minimum height=12mm, inner sep=3pt, font=\footnotesize},
    prod/.style={caixa, fill=gray!8},
    retre/.style={caixa, fill=gray!14},
    seta/.style={-{Stealth[length=2.6mm]}, thick},
    rot/.style={font=\scriptsize, align=center, inner sep=2pt}
]
    % ------------------ linha de cima: do laço à liberação -----------------
    \node[caixa] (laco) at (0,0)
        {\textbf{laço de aprendizado}\\\textbf{ativo}: seleção $\to$\\
         oráculo $\to$ retreino};
    \node[caixa] (parada) at (5.4,0)
        {\textbf{parada} pelo primeiro\\de três gatilhos:\\
         estagnação, orçamento\\ou exaustão do \textit{pool}};
    \node[caixa] (lib) at (10.8,0)
        {\textbf{liberação}: três condições;\\depois, uma única medida\\
         no teste intocado};

    % ------------------ linha de baixo: produção e retreino ----------------
    \node[prod] (producao) at (10.8,-3.6)
        {\textbf{produção}: monitora a\\deriva em três camadas\\
         (entrada; confiança;\\amostra com oráculo)};
    \node[retre, text width=44mm] (retreino) at (5.4,-3.6)
        {\textbf{retreino do zero} com\\
         \mbox{$L_{\text{antigo}} \cup L_{\text{novo}}$}, ciclo\\
         incremental no fluxo recente};

    % ------------------ setas do ciclo -------------------------------------
    \draw[seta] (laco) -- (parada);
    \draw[seta] (parada) -- (lib);
    \draw[seta] (lib) -- (producao);
    \draw[seta] (producao) -- (retreino);
    % o retorno fecha o ciclo pela esquerda, entrando por baixo do laço
    \draw[seta] (retreino.west) -| (laco.south);

    % ------------------ nota do ciclo de deriva ----------------------------
    \node[rot, text width=118mm] at (5.4,-5.6)
        {no monitoramento, as camadas 1 e 2 (custo zero) antecipam; a
         camada 3 (amostra com oráculo) confirma e dispara o retreino, que
         preserva todos os rótulos históricos e devolve o modelo à produção.};
\end{tikzpicture}
\end{document}
% GEOMETRIA CALIBRADA PARA CORPO 12 (ppginf/book 12pt, textwidth 16cm).
% Requer apenas arrows.meta e positioning (já em packages.tex). Legível em
% P&B: produção = cinza claro; retreino = cinza médio; operação = branco.
% ============================================================================
\begin{tikzpicture}[
    x=1cm, y=1cm,
    caixa/.style={draw, rounded corners=2pt, align=center, text width=40mm,
                  minimum height=12mm, inner sep=3pt, font=\footnotesize},
    prod/.style={caixa, fill=gray!8},
    retre/.style={caixa, fill=gray!14},
    seta/.style={-{Stealth[length=2.6mm]}, thick},
    rot/.style={font=\scriptsize, align=center, inner sep=2pt}
]
    % ------------------ linha de cima: do laço à liberação -----------------
    \node[caixa] (laco) at (0,0)
        {\textbf{laço de aprendizado}\\\textbf{ativo}: seleção $\to$\\
         oráculo $\to$ retreino};
    \node[caixa] (parada) at (5.4,0)
        {\textbf{parada} pelo primeiro\\de três gatilhos:\\
         estagnação, orçamento\\ou exaustão do \textit{pool}};
    \node[caixa] (lib) at (10.8,0)
        {\textbf{liberação}: três condições;\\depois, uma única medida\\
         no teste intocado};

    % ------------------ linha de baixo: produção e retreino ----------------
    \node[prod] (producao) at (10.8,-3.6)
        {\textbf{produção}: monitora a\\deriva em três camadas\\
         (entrada; confiança;\\amostra com oráculo)};
    \node[retre, text width=44mm] (retreino) at (5.4,-3.6)
        {\textbf{retreino do zero} com\\
         \mbox{$L_{\text{antigo}} \cup L_{\text{novo}}$}, ciclo\\
         incremental no fluxo recente};

    % ------------------ setas do ciclo -------------------------------------
    \draw[seta] (laco) -- (parada);
    \draw[seta] (parada) -- (lib);
    \draw[seta] (lib) -- (producao);
    \draw[seta] (producao) -- (retreino);
    % o retorno fecha o ciclo pela esquerda, entrando por baixo do laço
    \draw[seta] (retreino.west) -| (laco.south);

    % ------------------ nota do ciclo de deriva ----------------------------
    \node[rot, text width=118mm] at (5.4,-5.6)
        {no monitoramento, as camadas 1 e 2 (custo zero) antecipam; a
         camada 3 (amostra com oráculo) confirma e dispara o retreino, que
         preserva todos os rótulos históricos e devolve o modelo à produção.};
\end{tikzpicture}
\end{document}
% GEOMETRIA CALIBRADA PARA CORPO 12 (ppginf/book 12pt, textwidth 16cm).
% Requer apenas arrows.meta e positioning (já em packages.tex). Legível em
% P&B: produção = cinza claro; retreino = cinza médio; operação = branco.
% ============================================================================
\begin{tikzpicture}[
    x=1cm, y=1cm,
    caixa/.style={draw, rounded corners=2pt, align=center, text width=40mm,
                  minimum height=12mm, inner sep=3pt, font=\footnotesize},
    prod/.style={caixa, fill=gray!8},
    retre/.style={caixa, fill=gray!14},
    seta/.style={-{Stealth[length=2.6mm]}, thick},
    rot/.style={font=\scriptsize, align=center, inner sep=2pt}
]
    % ------------------ linha de cima: do laço à liberação -----------------
    \node[caixa] (laco) at (0,0)
        {\textbf{laço de aprendizado}\\\textbf{ativo}: seleção $\to$\\
         oráculo $\to$ retreino};
    \node[caixa] (parada) at (5.4,0)
        {\textbf{parada} pelo primeiro\\de três gatilhos:\\
         estagnação, orçamento\\ou exaustão do \textit{pool}};
    \node[caixa] (lib) at (10.8,0)
        {\textbf{liberação}: três condições;\\depois, uma única medida\\
         no teste intocado};

    % ------------------ linha de baixo: produção e retreino ----------------
    \node[prod] (producao) at (10.8,-3.6)
        {\textbf{produção}: monitora a\\deriva em três camadas\\
         (entrada; confiança;\\amostra com oráculo)};
    \node[retre, text width=44mm] (retreino) at (5.4,-3.6)
        {\textbf{retreino do zero} com\\
         \mbox{$L_{\text{antigo}} \cup L_{\text{novo}}$}, ciclo\\
         incremental no fluxo recente};

    % ------------------ setas do ciclo -------------------------------------
    \draw[seta] (laco) -- (parada);
    \draw[seta] (parada) -- (lib);
    \draw[seta] (lib) -- (producao);
    \draw[seta] (producao) -- (retreino);
    % o retorno fecha o ciclo pela esquerda, entrando por baixo do laço
    \draw[seta] (retreino.west) -| (laco.south);

    % ------------------ nota do ciclo de deriva ----------------------------
    \node[rot, text width=118mm] at (5.4,-5.6)
        {no monitoramento, as camadas 1 e 2 (custo zero) antecipam; a
         camada 3 (amostra com oráculo) confirma e dispara o retreino, que
         preserva todos os rótulos históricos e devolve o modelo à produção.};
\end{tikzpicture}
\end{document}
% GEOMETRIA CALIBRADA PARA CORPO 12 (ppginf/book 12pt, textwidth 16cm).
% Requer apenas arrows.meta e positioning (já em packages.tex). Legível em
% P&B: produção = cinza claro; retreino = cinza médio; operação = branco.
% ============================================================================
\begin{tikzpicture}[
    x=1cm, y=1cm,
    caixa/.style={draw, rounded corners=2pt, align=center, text width=40mm,
                  minimum height=12mm, inner sep=3pt, font=\footnotesize},
    prod/.style={caixa, fill=gray!8},
    retre/.style={caixa, fill=gray!14},
    seta/.style={-{Stealth[length=2.6mm]}, thick},
    rot/.style={font=\scriptsize, align=center, inner sep=2pt}
]
    % ------------------ linha de cima: do laço à liberação -----------------
    \node[caixa] (laco) at (0,0)
        {\textbf{laço de aprendizado}\\\textbf{ativo}: seleção $\to$\\
         oráculo $\to$ retreino};
    \node[caixa] (parada) at (5.4,0)
        {\textbf{parada} pelo primeiro\\de três gatilhos:\\
         estagnação, orçamento\\ou exaustão do \textit{pool}};
    \node[caixa] (lib) at (10.8,0)
        {\textbf{liberação}: três condições;\\depois, uma única medida\\
         no teste intocado};

    % ------------------ linha de baixo: produção e retreino ----------------
    \node[prod] (producao) at (10.8,-3.6)
        {\textbf{produção}: monitora a\\deriva em três camadas\\
         (entrada; confiança;\\amostra com oráculo)};
    \node[retre, text width=44mm] (retreino) at (5.4,-3.6)
        {\textbf{retreino do zero} com\\
         \mbox{$L_{\text{antigo}} \cup L_{\text{novo}}$}, ciclo\\
         incremental no fluxo recente};

    % ------------------ setas do ciclo -------------------------------------
    \draw[seta] (laco) -- (parada);
    \draw[seta] (parada) -- (lib);
    \draw[seta] (lib) -- (producao);
    \draw[seta] (producao) -- (retreino);
    % o retorno fecha o ciclo pela esquerda, entrando por baixo do laço
    \draw[seta] (retreino.west) -| (laco.south);

    % ------------------ nota do ciclo de deriva ----------------------------
    \node[rot, text width=118mm] at (5.4,-5.6)
        {no monitoramento, as camadas 1 e 2 (custo zero) antecipam; a
         camada 3 (amostra com oráculo) confirma e dispara o retreino, que
         preserva todos os rótulos históricos e devolve o modelo à produção.};
\end{tikzpicture}
